% !TeX root = ../main.tex
% 原始章节：07-exception.Rmd

\chapter{支持异常和中断}\label{chap:07-exceptions}

截至目前，本书中介绍的CPU设计仅处理了所有指令在没有错误的情况下，按照指令语义逐条执行的情形。给定一个程序（指令序列），计算机的执行过程是确定性的。然而，实际计算机工作过程中情况往往更加复杂。

本章将从现代计算机的并行计算和内存管理两个功能入手，讨论为什么CPU需要异常和中断机制，帮助读者理解异常和中断的功能。接下来，本章将结合LA32R架构，分析该架构中异常和中断的分类。然后，从CPU的角度详细介绍处理器对异常和中断的处理流程。最后，本章将面向CPU设计，针对具体的异常类别介绍设计要点。

\section{中断和异常机制和基本概念}\label{sec:07-exceptions-001}

\subsection{CPU的中断机制}\label{sec:07-exceptions-002}

并行计算是计算机实际使用中的一个基础内容，即使是单核CPU在操作系统的控制下也可以实现并行计算。操作系统通过分时机制，让需要并行执行的多个程序交替运行。例如，让程序A先执行若干个周期，再让程序B执行若干个周期。并行计算给程序员和用户提供了两种错觉：

\begin{enumerate}
\def\labelenumi{\arabic{enumi}.}
\item
  \textbf{程序员的错觉}：程序员只需独立编写程序A和程序B，并行计算无需修改这两个程序，只需通知操作系统即可。这就是现代并行计算给程序员的错觉，即在编程时无需考虑当前程序如何与其他程序并行执行。
\item
  \textbf{用户的错觉}：尽管只有一个CPU核心，操作系统通过将两个程序的时间片划分得足够小，用户会误以为两个程序在同时运行，这便是给用户的并行错觉。
\end{enumerate}

在本文已介绍的CPU架构中，如何为操作系统提供程序切换的支持？为了从程序A跳转到程序B，传统的方法可能需要在程序A中插入跳转或分支指令，并预先计算出程序B的入口地址，从而在执行到该指令时切换到程序B。这种方法需要修改应用程序，增加了开发的复杂性。

在实际的CPU中，实现并行计算效果依赖于CPU的\textbf{中断（Interruption）}机制。CPU中断机制是一种用于处理中断事件的硬件与软件相结合的技术。当外部设备、操作系统或程序自身发出中断请求时，CPU能够暂停当前正在执行的程序，保存其状态，并转而处理该中断请求。处理完中断后，CPU可以恢复先前程序的执行，从而保证程序执行的连续性。

操作系统利用中断机制实现并行计算。每当时间片结束或有重要事件发生时（例如，I/O完成或用户输入），操作系统会触发一个中断信号。CPU收到中断后，暂停当前程序的执行，保存其上下文信息，然后跳转到操作系统的中断处理程序。这一过程中，操作系统可以根据需要切换到另一个待执行的程序。通过这种方法，操作系统无需修改用户程序就能实现程序的快速切换，达到并行计算的效果。

\subsection{CPU的异常机制}\label{sec:07-exceptions-003}

内存管理是操作系统的重要功能之一，它确保多个程序能够在受限的物理内存中安全、高效地运行。为了让程序员能够在不考虑实际物理内存布局的情况下编写代码，操作系统引入了虚拟内存机制。虚拟内存为每个程序提供了独立的地址空间，使得程序员无需关心程序间的内存冲突问题。

在虚拟内存管理中，操作系统将物理内存划分为多个固定大小的页面，并将这些页面映射到程序的虚拟地址空间中。程序在执行时，通过CPU的内存管理单元（MMU）将虚拟地址转换为物理地址。这种地址转换使得程序可以在比实际物理内存更大的地址空间中运行，而操作系统负责管理这些映射关系。

然而，当程序访问的虚拟地址没有被映射到物理内存时，就会出现一个关键问题：CPU无法完成这次内存访问。在这种情况下，虚拟内存机制依赖于CPU提供的异常支持。当发生这种情况时，CPU会触发一个页错误异常，暂停当前程序的执行，并将控制权交给操作系统的异常处理程序。

操作系统通过异常处理程序可以加载所需的页面或处理其他内存管理问题，然后将控制权返回给CPU，继续程序的执行。通过这种方式，虚拟内存得以有效运作，而这一切都依赖于CPU的异常机制。

\textbf{异常（Exception）}是一种由程序执行过程中发生的错误或特殊事件触发的机制。当程序访问了无效的页面，或者发生了其他导致指令无法正常执行的情况时，CPU会自动暂停当前指令的执行，保存程序的状态，并跳转到异常处理程序。操作系统在异常处理程序中，可以采取多种措施来解决问题。

页错误异常指的是，当程序访问了一个未映射的虚拟地址时，CPU触发页错误异常。操作系统可以通过检查访问的地址，判断是否需要将对应的页面从磁盘加载到物理内存中，并更新页表。处理完这一异常后，操作系统可以让CPU恢复到错误发生前的状态，并重新执行引发异常的指令，从而让程序继续运行。

通过异常机制，操作系统能够灵活管理内存，提高系统的稳定性和健壮性。程序员不必在每个可能的内存访问错误上插入额外的错误检查逻辑，而是依赖于操作系统的异常处理机制，确保程序在遇到问题时能得到妥善处理。异常机制不仅使得程序更加简洁，还提高了系统资源的利用效率。

\subsection{基本概念}\label{sec:07-exceptions-004}

\textbf{异常（Exception）和中断（Interrupt）}是两类会打断CPU正常指令执行的事件。它们都使得CPU暂停当前的指令流，并转向预定义的处理程序，从而应对突发的情况。然而，二者在触发原因和处理方式上存在显著差异。

异常是由程序执行过程中发生的错误或特殊事件引发的。常见的异常包括除零错误、非法指令、页面错误（如访问未映射的虚拟地址）等。当CPU遇到这样的情况时，通常无法继续执行当前指令，必须进行特殊处理。此时，CPU会自动暂停当前程序的执行，保存当前的状态信息，然后跳转到操作系统定义的异常处理程序。异常处理程序可以修复错误、报告问题，或者在某些情况下终止问题程序。异常通常是同步事件，即它们与程序的执行紧密相关，并在指令执行时被检测到。

中断则是由外部事件或硬件设备引发的，例如键盘输入、硬盘I/O完成信号、定时器超时等。当这些设备或外部事件需要CPU立即响应时，它们会向CPU发送中断信号。收到中断信号后，CPU会暂停当前任务的执行，保存任务的状态，然后跳转到对应的中断处理程序。与异常不同，中断通常是异步事件，即它们与程序执行无关，可以在任意时间发生。中断的处理通常较为简短和快速，处理完后，CPU会恢复之前的任务继续执行。中断可以分为可屏蔽中断（maskable interrupts），如一般的外部设备中断，和不可屏蔽中断（non-maskable interrupts），如系统故障或其他紧急情况。

尽管异常和中断都涉及到CPU的状态切换与处理程序的执行，但它们的区别主要在于触发条件和性质。异常是由程序内部执行错误或特殊情况引发的同步事件，直接与当前指令流相关。而中断则是由外部硬件或其他系统事件引发的异步事件，与当前程序的指令流无关。此外，异常通常需要复杂的处理逻辑来修复程序状态或报错，而中断则通常只需处理简单的硬件请求，如读取数据或响应用户输入。

异常和中断在处理机制上具有许多相似之处。首先，二者都需要CPU跳转到指定的处理程序。其次，处理完毕后，CPU需要恢复之前的状态，继续执行被打断的程序。此外，操作系统的设计通常会统一管理异常和中断的处理机制，使用类似的框架来处理这两类事件。例如，异常和中断的处理程序都需要在系统的中断向量表或异常向量表中注册，并且都可能涉及对系统资源的调度和管理。通过统一的异常和中断处理机制，操作系统能够高效地应对各种突发事件，确保系统的稳定运行。

常见的异常事件包括：

\begin{itemize}
\tightlist
\item
  除零错误：当执行除法指令时，除数为零会导致异常，此时CPU通过异常机制进行处理。
\item
  非法指令：CPU遇到无法识别或无效的指令时触发异常。
\item
  页面错误：虚拟内存管理中，当访问一个未映射到物理内存的页面时会产生异常。
\end{itemize}

常见的中断事件包括：

\begin{itemize}
\tightlist
\item
  时钟中断：用于操作系统的调度程序切换任务。
\item
  外部中断：如键盘、鼠标或网络设备的输入信号，这些外设在工作完成或需要进行信息传递时，可通过使CPU陷入中断的方式通知CPU（还有其他的外设交互方法）。这种方式需要利用CPU的外部中断。
\item
  软件中断：软件可通过给体系结构规定的特定寄存器写入某些值的方式，使CPU陷入中断。与外部中断不同的是，这种中断时软件在运行时自行触发的。
\end{itemize}

\section{异常和中断的分类}\label{sec:07-exceptions-005}

\subsection{六类异常}\label{sec:07-exceptions-006}

在LoongArch架构中，异常可以根据其来源划分为以下六种类型，每种类型都有其独特的触发条件和处理要求：

1）\textbf{外部事件}：外部事件是指源自于CPU核心外部的模块或者通过实际的物理连线从处理器外部传入的事件，通常称为中断。中断使得CPU能够异步处理多个事件，尤其是在操作系统中，中断机制避免了轮询等待的资源浪费，特别是在处理与I/O相关的任务时显得尤为重要。中断事件的不可预测性以及无法由软件控制的特性，要求系统必须具备完善的软硬件协作机制，以确保中断的处理不会干扰程序的正常执行流程。例如，外部设备发出的I/O中断信号会打断当前执行的程序，转而执行中断服务程序（\texttt{ISR}），从而提高系统的响应速度。

2）\textbf{指令执行中的错误}：在指令执行过程中，可能会发生由于操作码或操作数不符合要求而引发的错误，例如调用不存在的指令、试图进行除以零的操作、访问不对齐的地址、在用户态下调用仅供核心态使用的指令或尝试访问非法的内存地址空间等。这些错误将导致指令无法继续执行，并且可能引发严重的系统错误。因此，系统必须具备相应的异常处理机制，以确保能够在发生这些错误时及时捕获并进行适当的处理，例如通过软件捕捉这些异常并采取相应的补救措施。

3）\textbf{数据完整性问题}：数据完整性问题通常涉及存储器数据的校验错误。例如，在存储器中使用ECC（错误检测和纠正码）等硬件校验方式时，如果发现数据在传输或存储过程中发生了校验错误，系统会引发异常。这些异常可以帮助系统识别出硬件中的潜在故障，从而采取措施，如重新读取数据或修复数据错误。对于可纠正的错误，系统可以自动修正；而对于不可纠正的错误，系统可能需要重新启动或切换到冗余系统以维持正常运行。

4）\textbf{地址转换异常}：在存储管理单元（\texttt{MMU}）尝试对内存页进行地址转换时，如果在硬件的页表或转换表中找不到相应的地址映射项，就会产生地址转换异常。这类异常通常表明程序试图访问的内存地址不存在或没有被正确映射，需要异常处理程序来解决这一问题。处理器可以通过异常处理程序将缺页错误（\texttt{Page\ Fault}）转交给操作系统，由操作系统决定是分配新的内存页，还是终止错误的进程。

5）\textbf{系统调用和陷入}：系统调用和陷入异常是由特定指令触发的，旨在生成操作系统可识别的信号，以便在受保护模式下调用核心态的相关操作。通常，这类异常用于执行需要更高权限的操作，例如文件系统访问、进程管理、硬件设备控制等。系统调用通过触发陷入异常，使得用户态程序能够安全地请求操作系统服务，而不必直接操作硬件资源，从而确保系统的安全性和稳定性。

6）\textbf{需要软件修正的运算}：在某些复杂的运算过程中，特别是涉及到浮点运算指令时，由于硬件实现的复杂性，可能无法直接处理所有情况。这时，处理器会引发异常，通知系统使用软件来完成这些复杂运算。例如，某些特殊的浮点运算或者未定义的操作会触发一个浮点异常，系统可以通过软件模拟的方式来完成这些计算，保证程序的正确性和稳定性。

\subsection{两类中断}\label{sec:07-exceptions-007}

在LoongArch架构中，中断是处理器应对外部事件的一种机制，通常通过两种主要方式传递到处理器：\textbf{中断线}和\textbf{消息中断}。

1）\textbf{中断线传递}：这是最直接的中断传递方式，通常用于中断源数量较少的系统。在这种方式中，中断源通过物理中断线直接连接到处理器的中断引脚。当处理器检测到某根中断线有信号时，处理器会优先处理该中断线所对应的中断事件，并调用相应设备的中断服务程序（\texttt{ISR}）。如果系统中有多个中断源，则可能通过中断控制器来汇总这些信号，进行中断优先级的管理和中断请求的分发。

2）\textbf{消息中断传递}：消息中断是一种更具扩展性和灵活性的方法，通过向特定的内存地址写入一个指定的数值来触发中断，而不依赖于物理中断线。这种方法不仅减少了对物理引脚的依赖，还允许系统动态地分配和管理中断源。消息中断的使用非常灵活，适用于复杂的系统设计，尤其是那些拥有大量设备和中断源的现代计算机系统。例如，在多核处理器或分布式系统中，消息中断可以用于实现高效的核间通信（\texttt{Inter-Processor\ Interrupts}，简称\texttt{IPI}），或者在虚拟化环境中，实现对虚拟机的中断传递和管理。

在LoongArch32R架构中，仅需要实现第一种中断机制。

\section{异常和中断处理流程}\label{sec:07-exceptions-008}

\subsection{异常处理流程}\label{sec:07-exceptions-009}

异常处理是系统处理异常事件的关键环节，通常包括以下五个阶段：\textbf{准备处理}、\textbf{确定异常来源}、\textbf{保存执行状态}、\textbf{处理异常}以及\textbf{恢复状态并返回}。这些阶段的核心目标是确保异常处理能够准确无误地进行，同时保持系统的整体稳定性。

1）\textbf{异常处理准备}：当异常发生时，CPU需要在进入异常处理程序之前完成一些准备工作。这包括记录异常发生时被打断的指令地址和处理状态。在LoongArch架构中，\texttt{CSR.ERA}寄存器用于保存发生异常时的指令地址，从而确保处理器能够在异常处理完成后，准确恢复到异常发生前的状态，继续执行中断的程序。此外，CPU通常会将其权限等级提升到最高特权级，并关闭中断响应。在LoongArch中，通过将\texttt{CSR.CRMD}寄存器的\texttt{PLV}域设置为0进入最高权限级别，并将\texttt{IE}域置0以屏蔽所有中断，确保异常处理过程中不被新的中断干扰。

在准备阶段，硬件还需保存当前的处理状态以供恢复。例如，LoongArch架构中，\texttt{CSR.CRMD}中的\texttt{PLV}和\texttt{IE}域的旧值会被存储到\texttt{CSR.PRMD}寄存器中，这样在异常处理结束后可以快速恢复之前的状态。同时，异常编号和其他相关信息会记录在\texttt{CSR.ESTAT}寄存器的相应域中，以供异常处理程序参考和处理。

2）\textbf{确定异常来源}：不同类型的异常需要特定的处理程序来应对。处理器通过异常编号来区分不同的异常类型，并跳转到相应的异常处理程序入口。在LoongArch架构中，异常处理程序入口地址由\texttt{CSR.EENTRY}寄存器和\texttt{CSR.ECFG}寄存器的组合来确定，每种异常类型都对应一个特定的入口地址。相比之下，在X86架构中，处理器使用中断描述符表（\texttt{IDT}）来确定异常处理程序的入口地址，每种异常类型都有其对应的中断描述符，描述符中包含了处理程序的段选择子和偏移地址。

3）\textbf{保存执行状态}：在进入异常处理程序之前，操作系统通常会将当前程序的状态保存到堆栈中。这包括通用寄存器、程序计数器（\texttt{PC}）、状态寄存器（\texttt{SR}）等关键数据的内容，以便在处理完异常后能够准确地恢复被打断的程序状态。这一步骤对于确保程序能够在异常处理完毕后继续正常执行至关重要。

4）\textbf{处理异常}：一旦异常处理程序被调用，系统会根据异常的类型执行特定的处理操作。例如，对于页面错误异常，处理程序可能需要分配新的页面并更新页表；对于系统调用异常，处理程序会根据调用的系统服务号执行相应的内核操作。在LoongArch中，异常处理程序可以通过查看\texttt{CSR.ESTAT}中的异常编号和相关信息来决定处理流程，并执行相应的指令以解决异常或完成所请求的操作。

5）\textbf{恢复执行状态并返回}：在异常处理完成后，系统需要将之前保存的执行状态恢复，并使程序继续从中断处或指定位置开始执行。具体的恢复操作包括从堆栈中取出保存的寄存器值、恢复程序计数器（\texttt{PC}）以及状态寄存器的内容。在LoongArch架构中，\texttt{CSR.ERA}寄存器的值会被恢复到\texttt{PC}中，确保程序能够从正确的位置继续执行。此外，\texttt{CSR.PRMD}寄存器中的旧权限等级和中断使能状态也会被恢复到\texttt{CSR.CRMD}中，以确保程序在处理完异常后回到正确的权限级别，并重新启用中断响应。

在一些复杂的异常处理中，异常处理程序可能需要进行额外的操作，如清理资源、记录日志或发出警告等。在此情况下，处理程序在返回之前，确保所有额外的操作都已经完成，且不会影响程序的正常运行。

恢复完成后，系统通过专门的指令返回到异常发生前的程序继续执行。例如，LoongArch使用\texttt{ERET}（Exception Return）指令来恢复异常前的状态并返回到中断的程序执行，这确保了系统在处理完异常后，能够准确无误地继续原有的程序流程。

\subsection{中断的响应}\label{sec:07-exceptions-010}

LoongArch指令系统共定义了13个中断，分别是：

\begin{enumerate}
\item 1个核间中断（\texttt{IPI}），
\item 1个定时器中断（\texttt{TI}），
\item 1个性能监测计数溢出中断（\texttt{PMI}），
\item 8个外部硬中断（\texttt{HWI0\textasciitilde{}HWI7}），
\item 2个软中断（\texttt{SWI0\textasciitilde{}SWI1}）。
\end{enumerate}


所有中断线上的中断信号均采用电平中断，并且都是高电平有效。当有中断发生时，这种高电平有效的中断方式会在输入到处理器的中断线上保持高电平状态，直到中断被处理器响应并处理为止。无论中断源来自处理器核的外部还是内部，是硬件触发还是软件置位，这些中断信号都会持续地被采样，并记录到\texttt{CSR.ESTAT}中的\texttt{IS}域对应的比特位上。

这些中断均为可屏蔽中断，除了\texttt{CSR.CRMD}中的全局中断使能位\texttt{IE}外，每个中断各自还有其局部中断使能控制位，位于\texttt{CSR.ECFG}的\texttt{LIE}域中。当\texttt{CSR.ESTAT}中的\texttt{IS}域的某位为1且对应的局部中断使能和全局中断使能均有效时，处理器将响应该中断，并进入中断处理程序的入口处开始执行。

在支持多个中断源输入的指令系统中，需要明确当多个中断同时触发时，处理器如何区别不同来源的中断优先级。

\textbf{非向量中断模式}：在非向量中断模式下不区分中断优先级。此时，如果需要对中断进行优先级处理，可以通过软件方式来实现。常见的实现方案包括：

\begin{enumerate}
\item 软件随时维护一个中断优先级（\texttt{IPL}），每个中断源都被赋予特定的优先级。
\item 在正常状态下，CPU运行在最低优先级，此时任何中断都可以触发。
\item 当CPU处于最高中断优先级时，任何中断都被禁止。
\item 当更高优先级的中断发生时，可以抢占低优先级的中断处理过程。
\end{enumerate}


\textbf{向量中断模式}：在向量中断模式下，处理器通常需要依照既定的优先级规则从多个已生效的中断源中选择一个，并跳转到其对应的处理程序入口处。完整LoongArch指令系统实现的是向量中断，采用固定优先级仲裁机制，具体规则如下：硬件中断号越大，优先级越高。也就是说，\texttt{IPI}的优先级最高，\texttt{TI}次之，依次递减，直到\texttt{SWI0}，其优先级最低。由于简化版的LA32R架构无需实现向量中断，因此LA32R的开发者也不需要处理向量中断中的优先级问题。

\section{优先级}\label{sec:07-exceptions-011}

在处理器设计中，异常和中断是两类不同的事件处理机制，它们各自承担着不同的功能。\textbf{异常}通常是由于程序执行中的错误或特定指令而触发的，例如除以零、非法指令或地址不对齐等情况。而\textbf{中断}则是由外部事件或硬件信号触发的，如定时器中断、外设中断等，用于打断当前程序的执行，以便处理紧急任务。

\subsection{异常与中断的优先级}\label{sec:07-exceptions-012}

在大多数处理器架构中，\textbf{异常}的优先级通常高于\textbf{中断}。这是因为异常通常与当前指令的执行密切相关，并且如果不立即处理，可能会导致系统的不稳定或崩溃。例如，当遇到非法指令或除零错误时，处理器必须立即响应并处理，否则程序的执行将无法继续。

另一方面，中断尽管需要及时处理，但通常它们与外部事件有关，通常可以在异常处理完成后再进行响应。这种设计确保了系统在遇到严重错误时能够优先处理这些错误，维持系统的稳定性。

\subsection{精确异常}\label{sec:07-exceptions-013}

\textbf{精确异常}是指在发生异常时，处理器能够保证异常发生前的所有指令已经正确执行，而异常发生后的指令未执行或可以撤销。这意味着，当异常处理完成后，处理器能够从异常发生的指令重新开始执行，而不会丢失或重复任何指令的执行。

精确异常的实现对于处理器的可靠性至关重要，尤其是在复杂的流水线和超标量处理器中。在这些处理器中，指令的执行顺序可能与程序的顺序不同，因此在异常发生时，确保异常精确定位和处理成为了处理器设计中的一个重要挑战。

在大多数现代处理器中，硬件支持精确异常的机制。例如，在异常发生时，处理器会记录被打断的指令地址（通常存储在如\texttt{CSR.ERA}的寄存器中），确保异常处理完成后能够正确地恢复执行。

\subsection{中断的异步响应}\label{sec:07-exceptions-014}

\textbf{中断}的一个关键特性是其\textbf{异步性}。与异常不同，中断并不与指令的执行直接相关。中断可以在任何时候发生，而不受当前正在执行的指令的限制。这种异步性使得中断成为处理外部事件的有效机制，例如定时器溢出或外设的输入信号。

当中断发生时，处理器会暂停当前的指令执行，保存上下文，然后跳转到中断处理程序。在LoongArch架构中，中断信号被不断采样并记录在\texttt{CSR.ESTAT}的\texttt{IS}域中。处理器根据中断优先级决定是否响应中断，并通过跳转到中断处理程序的入口地址来处理该中断。

在实际的CPU开发过程中，中断的异步响应特性给测试工作带来了很大的困难。由于不同微架构的处理器可能对于同时出现的中断信号有不同的响应时间，因此往往没有一个固定的正确行为可供开发人员比对。因此，开发者要么对中断进行手动测试，要么开发或使用一些测试框架。本团队的敏捷开发测试框架可缓解异步中断带来的测试困难，具体可参考另一个文档。

\subsection{中断和异常的嵌套处理}\label{sec:07-exceptions-015}

在复杂的计算系统中，处理器可能会在处理一个异常或中断时再次遇到新的异常或中断，这就是中断和异常的嵌套处理。为了处理这种情况，LoongArch架构提供了灵活的机制来管理和处理嵌套的中断和异常。

1）\textbf{优先级机制}：嵌套处理中，处理器需要通过优先级机制来决定是立即处理新的异常/中断，还是先完成当前的异常/中断处理。不同的异常和中断有不同的优先级，通常外部中断的优先级最低，而系统错误类异常的优先级最高。例如，在LoongArch中，当处理一个低优先级的中断时，如果一个更高优先级的异常发生，处理器会暂时中断当前的中断处理，转而处理高优先级的异常。处理完高优先级异常后，处理器会恢复到之前的中断处理流程。

2）\textbf{状态保存与恢复}：为了正确处理嵌套的中断和异常，处理器在每次进入新的异常或中断处理程序时，都需要保存当前的状态，并在退出时恢复之前的状态。这通常涉及堆栈的使用，将每一层嵌套的上下文信息（如寄存器内容、程序计数器等）压入堆栈，以便在处理完当前的异常/中断后能够恢复到正确的处理状态。例如，在LoongArch中，每次嵌套时，\texttt{CSR.ERA}、\texttt{CSR.CRMD}、\texttt{CSR.PRMD}等寄存器的值都会保存到堆栈中，确保嵌套处理不会丢失原有的处理状态。

3）\textbf{中断屏蔽与使能}：在处理嵌套的异常/中断时，处理器可以通过屏蔽某些中断来防止进一步的嵌套。例如，当处理高优先级的异常时，可以暂时屏蔽低优先级的中断，确保处理器不被其他不重要的中断打扰。LoongArch中的\texttt{CSR.CRMD}寄存器的\texttt{IE}位用于控制全局中断使能，通过设置或清除这个位，处理器可以有效管理在嵌套处理期间的中断行为。

4）\textbf{返回机制}：在嵌套处理完成后，处理器需要逐层恢复之前的中断或异常处理状态，并正确返回到中断或异常发生前的程序执行。在LoongArch架构中，使用\texttt{ERET}指令不仅用于简单的异常返回，还能处理嵌套的情况，通过\texttt{CSR.PRMD}和\texttt{CSR.ERA}寄存器的内容，处理器可以正确返回到之前的异常或中断处理程序中，逐层退出嵌套，直至最终返回到最初的程序执行。

在LA32R中不需要支持中断的嵌套处理机制，但对嵌套机制的了解有助于进行CPU的设计，以及对软硬件协同工作原理的理解。

\section{在CPU中增加异常与中断功能的设计要点}\label{sec:07-exceptions-016}

在CPU设计中，增加异常和中断功能是实现系统健壮性和灵活性的关键步骤。为此，开发者需要从多个方面考虑，确保这些功能能够可靠且高效地处理各种异常和中断情况。

在处理器设计中，几种常见的异常和中断需要特别关注，包括\textbf{地址错}、\textbf{整形溢出}以及\textbf{系统调用}。

\begin{itemize}
\item
  \textbf{地址错异常}通常由无效的内存访问引发，比如访问未对齐的地址或越界访问。这类异常的检测逻辑通常在地址生成单元（AGU）中实现，通过检查访问地址是否符合对齐要求，以及是否在合法的地址空间范围内。如果检测到地址错误，处理器会触发异常，并跳转到对应的异常处理程序。
\item
  \textbf{整形溢出异常}发生在算术运算结果超出数据类型的表示范围时。例如，在32位整数加法中，如果结果超出32位表示范围，处理器会产生溢出异常。为了检测这种情况，算术逻辑单元（ALU）会在执行运算后检查结果的符号位与溢出位，确认是否发生了溢出。
\item
  \textbf{系统调用}是一种特殊的异常，用于在用户态下调用操作系统内核功能。系统调用通常由专用指令（如LoongArch中的\texttt{SYSCALL}指令）触发。系统调用异常的检测逻辑相对简单，只需检测到指令后触发相应的异常处理即可。
\end{itemize}

\subsection{精确异常与中断的实现}\label{sec:07-exceptions-017}

根据上文的介绍，\textbf{精确异常}是指在异常发生时，确保被异常打断的指令及其之前的指令已经完整执行，而之后的指令未被执行。在具体进行开发实现时，开发人员需要在硬件和控制逻辑中精确跟踪指令的执行状态。当发生异常时，开发人员应让处理器保存当前指令的地址（如\texttt{CSR.ERA}寄存器），并正确地维护\texttt{CSR.CRMD}和\texttt{CSR.ESTAT}等寄存器的状态，以便异常处理程序能够准确恢复并继续执行。

\begin{itemize}
\tightlist
\item
  \textbf{\texttt{CSR.ERA}}寄存器保存异常发生时的指令地址，允许异常处理结束后恢复执行。
\item
  \textbf{\texttt{CSR.CRMD}}寄存器控制处理器的模式和全局中断使能，异常发生时会自动调整权限等级。
\item
  \textbf{\texttt{CSR.ESTAT}}寄存器记录当前异常或中断的状态，包括异常编号等关键信息。
\end{itemize}

\subsection{添加异常中断的支持}\label{sec:07-exceptions-018}

要在CPU中添加系统调用异常中断的支持，首先需要在指令解码阶段识别系统调用指令。当指令解码单元识别到\texttt{SYSCALL}指令时，它会触发异常处理流程。处理器需要跳转到系统调用异常处理程序的入口地址，并保存相关上下文信息（如\texttt{CSR.ERA}和\texttt{CSR.CRMD}）。在异常处理程序中，操作系统会根据系统调用号执行相应的内核功能，并最终通过执行异常返回指令（如LoongArch中的\texttt{ERTN}指令）恢复用户态的执行。

以\textbf{地址错异常}为例，添加支持的过程类似于系统调用异常。在地址生成单元中，检测到地址错时，处理器将生成一个地址错异常信号，并保存导致异常的虚拟地址到特定寄存器（如\texttt{CSR.BADV}）。接下来，处理器会跳转到地址错异常处理程序，该程序可以记录错误信息、通知操作系统或采取修复措施。

通过逐步实现和集成这些异常和中断的检测逻辑和处理流程，开发者可以确保处理器能够正确处理各种异常情况，提高系统的可靠性和稳定性。为了确保正确性，开发者应仔细阅读指令集手册，并充分理解处理器架构中的各类寄存器和控制状态。
